\section{Introduction}\label{sy-0100-ai}

This paper studies widgets. The basic identity is
\begin{equation}\label{eq:prose-ai}
\sigma^2 = \mathrm{id},
\end{equation}
which we use throughout.\footnote{The name comes from nowhere.} The \GW{} theory of gadgets is not touched. See \cite[Chapter II]{Har77} for schemes and \cite[Theorem 2.1]{Kre99} for cycle groups; the lemma \cite[Lemma 9]{Kre99} is not in the digest.

\begin{definition}[Widget]\label{sy-0001-ai}
A \emph{widget} is a pair $(X,\sigma)$ of a set and an involution $\sigma$ of $X$, so that \eqref{eq:prose-ai} holds. Its \emph{fixed locus} is
\begin{equation}\label{eq:fix-ai}
\Fix(\sigma) = \{x \in X : \sigma x = x\}.
\end{equation}
\end{definition}

% !LOOM author: The synthetic quilt
% !LOOM created: 2026-09-16
% !LOOM tags: orbits, setup

\begin{remark}\label{sy-000A-ai}
Gadgets exist: take $X = \{1,2\}$ with the swap. Compare \cite[Theorem 2.1]{Kre99}.
\end{remark}

\begin{lemma}[Orbits]\label{sy-0002-ai}
Every orbit of a widget has at most two points, and one exactly when its point is fixed; compare \eqref{eq:fix-ai} in Definition~\ref{sy-0001-ai}.
\end{lemma}
\begin{proof}
The orbit of $x$ is $\{x, \sigma x\}$ because
\begin{equation}\label{eq:step-ai}
\sigma(\sigma x) = x
\end{equation}
by \eqref{eq:prose-ai}.
\end{proof}


\begin{itemize}
\item widgets,
\item gadgets.
\end{itemize}

\begin{tabular}{ll}
object & involution \\
widget & yes \\
gadget & no
\end{tabular}

\begin{center}
\begin{tikzcd}
X \arrow[r, "\sigma"] \arrow[d] & X \arrow[d] \\
Y \arrow[r] & Y
\end{tikzcd}
\end{center}

\includegraphics[width=3cm]{figures/fig.pdf}

\parbox{5cm}{A boxed paragraph the converter renders as an image.}

\red{This sentence is in red.}

% !LOOM author: The synthetic quilt
% !LOOM tags: results

\begin{definition}[Orbit pair]\label{sy-999C-ai}
An \emph{orbit pair} is an orbit with two points.
\end{definition}

